% =====================================================================
%  Material posterior — Fuentes, bibliografía y licencia
%  La lista de referencias se genera con BibTeX a partir de
%  bibliography/referencias.bib (estilo autor-año, natbib + plainnat).
% =====================================================================
\chapter{Fuentes y bibliografía}

Fuentes utilizadas para la actualización de esta edición. Se ha procurado emplear fuentes oficiales y publicaciones académicas; el texto del manual es una elaboración original y no reproduce literalmente dichas obras.

% Fuentes citadas en prosa o consultadas que deben figurar en la lista
% aunque no lleven una cita \citep explícita en el cuerpo del texto:
\nocite{codigopenal1995, alertcops, igualdad, navarro}

% Suprimimos el encabezado propio de natbib: el título de capítulo ya
% lo proporciona el \chapter de arriba.
\renewcommand{\bibsection}{}
\bibliography{bibliography/referencias}

\section*{Licencia}
\addcontentsline{toc}{section}{Licencia}

Esta obra está bajo una licencia Reconocimiento-No comercial-Compartir bajo la misma licencia 3.0 Unported (CC BY-NC-SA 3.0) de Creative Commons. Para ver una copia de esta licencia, visite \url{http://creativecommons.org/licenses/by-nc-sa/3.0/} o envíe una carta a Creative Commons, 171 Second Street, Suite 300, San Francisco, California 94105, USA. Las obras derivadas deben distribuirse bajo esta misma licencia.

Si estás interesada en asistir al próximo seminario de Defensa Personal Femenina, visita \url{https://www.artesmarcialesbilbao.com}. Si tienes un grupo y deseas una convocatoria extraordinaria, contacta con \href{mailto:admin@artesmarcialesbilbao.com}{admin@artesmarcialesbilbao.com}.
